% GENERATED FILE. Edit canonical lessons, then run npm run generate:books.
% canonical-source-hash: fnv1a64:0acfaa551d1665f4
% canonical-lessons: IT-C114-lezione, IT-C114-esame, IT-C114-parola, IT-C114-lavagna, IT-C114-gomma

\chapter{Lesson, Exam, Word, Board, Rubber}
\label{ch:it-lesson-exam-word-board-rubber}
\hlchaptermodality{114}

\begin{chapteropening}
\textbf{By the end of this chapter:} \emph{I can say a lesson, an exam, a word, a board and a rubber.}

Complete the last lesson of chapter 114: I can say a lesson, an exam, a word, a board and a rubber.
\end{chapteropening}

\section[la lezione]{la lezione — a lesson}

\label{lesson:IT-C114-lezione}



\begin{quote}
Before the new one: say the Italian for a till, then the Italian for a class.

\emph{Read it:} \emph{il cinema, il teatro, lo stadio} aloud as one line, without stopping between the words
\end{quote}

\subsection*{You'll want to know: la lezione}
\textbf{la lezione} — \textquotedblleft{}a lesson\textquotedblright{}.

\begin{cousinweb}
\textbf{lezione} $\leftarrow$ Latin \emph{lectio}, \textquotedblleft{}a reading\textquotedblright{}, from \emph{legere}, \textquotedblleft{}to read\textquotedblright{} — a learned borrowing. (EtymDB 2.1)
\end{cousinweb}

\begin{grammarlens}[title={What you've built}]
One more word for this chapter.
\end{grammarlens}

\subsection*{Your turn}
\begin{itemize}
\raggedright
  \item \emph{Say it:} \emph{la lezione}
  \item \emph{Say it:} \emph{la lezione}, once more
  \item \emph{Say it:} say \emph{la lezione}
  \item \emph{Recall:} say the Italian for a till, then the Italian for a class, then say \emph{la lezione} again
\end{itemize}

\begin{tcolorbox}[breakable,colback=teal!4,colframe=teal!35!black,title={Before you move on}]
What does la lezione mean? (\textquotedblleft{}A lesson\textquotedblright{}.) Say it once more.
\end{tcolorbox}
\section[l'esame]{l'esame — an exam}

\label{lesson:IT-C114-esame}



\begin{quote}
Before the new one: say the Italian for a class, then the Italian for a lesson.
\end{quote}

\subsection*{You'll want to know: l'esame}
\textbf{l'esame} — \textquotedblleft{}an exam\textquotedblright{}.

\begin{cousinweb}
\textbf{esame} $\leftarrow$ Latin \emph{examen}, \textquotedblleft{}the tongue of a balance, a weighing\textquotedblright{} (Lewis \& Short) — a learned borrowing. (EtymDB 2.1)
\end{cousinweb}

\begin{grammarlens}[title={What you've built}]
One more word for this chapter.
\end{grammarlens}

\subsection*{Your turn}
\begin{itemize}
\raggedright
  \item \emph{Say it:} \emph{l'esame}
  \item \emph{Say it:} \emph{l'esame}, once more
  \item \emph{Say it:} say \emph{l'esame}
  \item \emph{Recall:} say the Italian for a class, then the Italian for a lesson, then say \emph{l'esame} again
\end{itemize}

\begin{tcolorbox}[breakable,colback=teal!4,colframe=teal!35!black,title={Before you move on}]
What does l'esame mean? (\textquotedblleft{}An exam\textquotedblright{}.) Say it once more.
\end{tcolorbox}
\section[la parola]{la parola — a word}

\label{lesson:IT-C114-parola}



\begin{quote}
Before the new one: say the Italian for a lesson, then the Italian for an exam.
\end{quote}

\subsection*{You'll want to know: la parola}
\textbf{la parola} — \textquotedblleft{}a word\textquotedblright{}.

\begin{cousinweb}
\textbf{parola} $\leftarrow$ Late Latin \emph{parabola}, \textquotedblleft{}a word, speech\textquotedblright{} — inherited. (EtymDB 2.1) In older Latin \emph{parabola} meant \textquotedblleft{}a comparison\textquotedblright{}, taken from Greek \emph{parabole}. (Lewis \& Short)
\end{cousinweb}

\begin{grammarlens}[title={What you've built}]
One more word for this chapter.
\end{grammarlens}

\subsection*{Your turn}
\begin{itemize}
\raggedright
  \item \emph{Say it:} \emph{la parola}
  \item \emph{Say it:} \emph{la parola}, once more
  \item \emph{Say it:} say \emph{la parola}
  \item \emph{Recall:} say the Italian for a lesson, then the Italian for an exam, then say \emph{la parola} again
\end{itemize}

\begin{tcolorbox}[breakable,colback=teal!4,colframe=teal!35!black,title={Before you move on}]
What does la parola mean? (\textquotedblleft{}A word\textquotedblright{}.) Say it once more.
\end{tcolorbox}
\section[la lavagna]{la lavagna — a board}

\label{lesson:IT-C114-lavagna}



\begin{quote}
Before the new one: say the Italian for an exam, then the Italian for a word.
\end{quote}

\subsection*{You'll want to know: la lavagna}
\textbf{la lavagna} — \textquotedblleft{}a board\textquotedblright{}.

\begin{cousinweb}
\textbf{lavagna} is named after Lavagna, a town on the Ligurian coast whose slate was cut into writing boards. EtymDB 2.1 has no row for it.
\end{cousinweb}

\begin{grammarlens}[title={What you've built}]
One more word for this chapter.
\end{grammarlens}

\subsection*{Your turn}
\begin{itemize}
\raggedright
  \item \emph{Say it:} \emph{la lavagna}
  \item \emph{Say it:} \emph{la lavagna}, once more
  \item \emph{Say it:} say \emph{la lavagna}
  \item \emph{Recall:} say the Italian for an exam, then the Italian for a word, then say \emph{la lavagna} again
\end{itemize}

\begin{tcolorbox}[breakable,colback=teal!4,colframe=teal!35!black,title={Before you move on}]
What does la lavagna mean? (\textquotedblleft{}A board\textquotedblright{}.) Say it once more.
\end{tcolorbox}
\section[la gomma]{la gomma — a rubber, an eraser}

\label{lesson:IT-C114-gomma}



\begin{quote}
Before the new one: say the Italian for a word, then the Italian for a board.
\end{quote}

\subsection*{You'll want to know: la gomma}
\textbf{la gomma} — \textquotedblleft{}a rubber, an eraser\textquotedblright{}.

\begin{cousinweb}
\textbf{gomma} $\leftarrow$ Late Latin \emph{gumma}, from Latin \emph{gummi}, \textquotedblleft{}gum\textquotedblright{} — inherited; Latin took it from Greek \emph{kommi}. (EtymDB 2.1)
\end{cousinweb}

\begin{grammarlens}[title={What you've built}]
One more word for this chapter.
\end{grammarlens}

\subsection*{Your turn}
\begin{itemize}
\raggedright
  \item \emph{Say it:} \emph{la gomma}
  \item \emph{Say it:} \emph{la gomma}, once more
  \item \emph{Say it:} say \emph{la gomma}
  \item \emph{Recall:} say the Italian for a word, then the Italian for a board, then say \emph{la gomma} again
\end{itemize}

\begin{tcolorbox}[breakable,colback=teal!4,colframe=teal!35!black,title={Before you move on}]
What does la gomma mean? (\textquotedblleft{}A rubber, an eraser\textquotedblright{}.) Say it once more.
\end{tcolorbox}
